We ran VERA with its frozen InternVL2-8B backbone on all \val{videos} videos of the VISILAB firearm-action dataset (\val{windows} ten-second windows) in three conditions on identical eight-frame inputs: VERA unchanged (A), firearm grounding only (G), and VERA with a grounding request in the same answer (AG). Unchanged VERA labels \val{T3_All_anomalypct}\,\% of windows anomalous, including \val{T3_NoGun_anomalypct}\,\% of windows without a firearm. Under VERA's packing of eight frames into one image block, the backbone does not localize firearms: G and AG leave almost every slot empty, every box they write covers the whole frame, and recall against the \val{T8_primary_G_All_iou50_gtobjects} annotated firearms is zero at IoU 0.50 and 0.75. Requesting localization still changes the anomaly judgments: AG flags \val{T4_All_diffpp}\,pp fewer windows than A (95\,\% CI \val{T4_All_diffppcilow} to \val{T4_All_diffppcihigh}), most of all on no-firearm windows. Because AG produced no localization, this is the effect of a bundled prompt change, not of localization. The study is exploratory: the grounding prompt's format was fixed after a bounded pilot.
